% ==================================================================
% id:       der-009
% topic:    נגזרות
% subtopic: הגדרת הנגזרת וגזירות
% solution: full
% ==================================================================
נתבונן בפונקציה $\ds f(x)=\case{cc}{x^3+2x^2-5x+2 & x<1 \\ x^4-2x^3+4x & x\geq 1}$. האם $f(x)$ גזירה ב $x=1$?

\ifsol
\hint{בדקו בנקודת התפר את כל התנאים הדרושים לגזירות, לא רק את הנגזרות של שני הענפים.}
\begin{solbox}
כמו בשאלות הקודמות, נסמן
$g(x)=x^3+2x^2-5x+2$,
$h(x)=x^4-2x^3+4x$.
מתקיים
$$,f'_{-}(1)=g'_{-}(1)=2$$
כמוכן
$$.f'_{+}(1)=h'_{+}(1)=2$$
אנו רואים כי
$$.f'_{-}(1)=f'_{+}(1)$$
אך האם זה מספיק לקביעת הגזירות של
$f$
ב-
$x=x_0$?
ממש לא!
נשים לב כי
$f$
אפילו לא רציפה, לכן, אל אחת כמה וכמה, לא גזירה.
ואמנם
    $$,\lim_{x\to 1^-}f(x)=\lim_{x\to 1^-}(x^3+2x^2-5x+2)=0$$
מצד שני
$$.\lim_{x\to 1^+}f(x)=\lim_{x\to 1^+}(x^4-2x^3+4x)=3$$
\end{solbox}
\fi
